\documentclass[12pt]{bxjsarticle}
\usepackage{amsmath,amssymb,bm}
\usepackage{booktabs}
\usepackage{geometry}
\geometry{margin=25mm}
\setlength{\parskip}{0.6em}
\setlength{\parindent}{0pt}
\title{2x2行列による動的システムの挙動解析：固有値と対角化の応用}
\date{}
\author{}
\begin{document}
\maketitle

\section{線形写像としての行列とシステム挙動の基礎}

動的システムの解析において、行列を単なる数値の配列としてではなく、ベクトルを別のベクトルへと変換する「線形写像」として捉えることは、物理的・工学的本質を理解するための戦略的必須条件である。システムの状態が時間とともに遷移する際、この写像は「現在の状態」から「次ステップの状態」への決定論的な推移を規定する。

\subsection{ベクトルの線形結合と座標表現の再定義}

実務的な状態空間表現（State-space representation）において、任意のシステム状態ベクトル  $\mathbf{u}$  は、基底ベクトル  $\mathbf{e}_1, \mathbf{e}_2$  の線形結合として定義される。  $$\mathbf{u} = u_1\mathbf{e}_1 + u_2\mathbf{e}_2$$  ここで、 $\mathbf{e}_1 = \begin{pmatrix} 1 \\ 0 \end{pmatrix}, \mathbf{e}_2 = \begin{pmatrix} 0 \\ 1 \end{pmatrix}$  とすれば、システムは直交座標系における成分として記述される。

\subsection{行列  $A$  による変換のメカニズム}

行列  $A = (a_{ij})$  がベクトル  $\mathbf{u}$  に作用するプロセスは、基底の変換として解釈できる。線形性の定義  $A(k\mathbf{u} + l\mathbf{v}) = kA\mathbf{u} + lA\mathbf{v}$  に基づき、写像後のベクトル  $\mathbf{v} = A\mathbf{u}$  は以下のように展開される。  $$\mathbf{v} = A(u_1\mathbf{e}_1 + u_2\mathbf{e}_2) = u_1(A\mathbf{e}_1) + u_2(A\mathbf{e}_2)$$  すなわち、行列  $A$  の各列ベクトルは「変換後の新しい基底」を示しており、この配置がシステム空間の伸縮や回転を決定づける。

\subsection{離散時間動的システムへの適用}

状態推移を  $\mathbf{u}_{n+1} = A\mathbf{u}_n$  と定義すると、初期状態  $\mathbf{u}_0$  から  $n$  ステップ後の状態は  $\mathbf{u}_n = A^n \mathbf{u}_0$  となる。この反復演算がシステムの漸近挙動（Asymptotic behavior）を支配するため、行列  $A$  の累乗特性を解析することがシステム予測の核心となる。

\section{固有方程式とケーリー・ハミルトンの定理による特性評価}

システムの運命を決定づけるのは、変換において「方向を変えず、長さのみが変化する」特別なベクトル、すなわち固有ベクトルとそれに対応する固有値  $\lambda$  である。

\subsection{固有方程式の導出と不変量}

固有値  $\lambda$  を特定するための特性方程式は、 $\det(A - \lambda I) = 0$  より導出される。行列  $A = \begin{pmatrix} a & b \\ c & d \end{pmatrix}$  において、これは以下の2次方程式に帰着する。  $$\lambda^2 - \text{tr}(A)\lambda + \det A = 0$$  ここで、トレース  $\text{tr}(A) = a + d$  および行列式  $\det A = ad - bc$  は、基底の取り方に依存しない行列の不変量（インバリアント）であり、システムの物理的性質を規定する。

\subsection{ケーリー・ハミルトンの定理}

システム解析における重要なマイルストーンとして、ケーリー・ハミルトンの定理が挙げられる。  $$A^2 - \text{tr}(A)A + (\det A)I = O$$  この定理は、行列自身がその特性方程式を満たすことを示しており、高次の累乗  $A^n$  を低次の行列の線形結合に簡約化する際の数学的根拠となる。

\subsection{判別式による挙動予測}

判別式  $D = (\text{tr}(A))^2 - 4\det A$  の符号は、システムの安定性と推移パターンを3つのシナリオに分類する。

\begin{center}
\begin{tabular}{lll}
\toprule
判別式 & 固有値の性質 & システム特性 \\
\midrule
$D > 0$ & 相異なる2つの実数解 & 2方向への独立した指数的伸縮 \\
$D = 0$ & 実数の重解 & 臨界状態、または線形成長を伴う推移 \\
$D < 0$ & 共役な複素数解 & 回転および振動成分の介在 \\
\bottomrule
\end{tabular}
\end{center}

\section{固有値のパターン別挙動解析と「標準形（Standard Forms）」}

システム設計の最適化には、固有値の構成に応じた「標準形」の理解が不可欠である。これにより、複雑な座標系での挙動を最も単純な構造に還元できる。

\subsection{相異なる2つの実数固有値 ( $D > 0$ )：対角形}

2つの独立した固有ベクトルが存在し、行列は以下の \textbf{標準形} へと対角化される。  $$\begin{pmatrix} \lambda_1 & 0 \\ 0 & \lambda_2 \end{pmatrix}$$  システムは各固有ベクトル方向へ個別に  $n$  乗の速度で成長または減衰する。干渉のない純粋な伸縮挙動である。

\subsection{複素固有値 ( $D < 0$ )：回転行列}

固有値が  $\mu \pm i\nu$  の場合、システムには回転成分が現れる。実数範囲での解析において、以下の \textbf{標準形} が定義される。  $$\begin{pmatrix} \mu & \nu \\ -\nu & \mu \end{pmatrix}$$  ここで、 $\mu$  は「利得（Gain）/減衰（Decay）」を制御し、 $\nu$  は「角速度（Angular Velocity）」に対応する。システムは複素平面上（あるいは実平面での射影）で螺旋を描きながら推移する。

\subsection{重解 ( $D = 0$ )：ジョルダン標準形と非対称推移}

固有値が  $\lambda$  の重解となる場合、2つのケースに分かれる。

\begin{enumerate}
\item \textbf{スカラ行列:}   $b=c=0$  の場合、 $A = \lambda I$ 。全方向へ均等に伸縮する。
\item \textbf{非対角化可能ケース:}  適切な基底変換により、以下の \textbf{標準形} が得られる。  $$\begin{pmatrix} \lambda & 1 \\ 0 & \lambda \end{pmatrix}$$  この構造では、指数的な伸縮  $\lambda^n$  に加え、非対角成分の「1」に起因する \textbf{線形成長} （ $n \lambda^{n-1}$  の項）が重畳される。これはシステム内に非対称な干渉が発生していることを意味する。
\end{enumerate}

\section{対角化を用いたシステム推移の高速演算と安定性評価}

大規模な動的システムにおいて、長期的な挙動を  $O(n)$  のステップ計算で追うのは非効率である。対角化は、計算コストを  $O(1)$  へと削減する戦略的な武器となる。

\subsection{計算効率の最大化}

固有ベクトルを並べた正則行列  $P$  を用いて  $A = PDP^{-1}$  と分解すると、累乗計算は中間項が相殺され、中心の対角行列  $D$  の累乗のみに帰着する。  $$A^n = P D^n P^{-1} = P \begin{pmatrix} \lambda_1^n & 0 \\ 0 & \lambda_2^n \end{pmatrix} P^{-1}$$  直接の行列積を  $n$  回繰り返す計算コストを回避し、任意のタイムステップ  $n$  における状態を瞬時に算定可能にする。

\subsection{システム安定性の判定：スペクトル半径  $\rho(A)$}

システムの長期的な「漸近安定性（Asymptotic Stability）」は、固有値の最大絶対値、すなわちスペクトル半径  $\rho(A) = \max |\lambda_i|$  によって決定される。

\begin{enumerate}
\item \textbf{$\rho(A) < 1}$  \textbf{:}  全ての固有値の絶対値が1未満。システムは  $n \to \infty$  で原点へと収束する（安定）。
\item \textbf{$\rho(A) > 1}$  \textbf{:}  少なくとも1つの固有値が単位円の外側にある。システムは発散する（不安定）。
\item \textbf{$\rho(A) = 1}$  \textbf{:}  境界状態。複素固有値の場合は永続的な振動（リミットサイクル）へと繋がる可能性がある。
\end{enumerate}

\section{結論：線形代数によるシステム設計の最適化}

本解析を通じて明らかなように、2x2行列の固有値解析は、単なる計算手法ではなく、システムの動態を規定する「普遍的な言語」である。

\begin{enumerate}
\item \textbf{定量的予測の実現:}  トレース、行列式、およびスペクトル半径を監視することで、シミュレーションを回す前にシステムの安定限界を定量的に評価できる。
\item \textbf{設計の合理化:}  「標準形」への還元は、複雑な干渉を排した理想的な座標系での設計を可能にし、制御戦略の構築において最短経路を提供する。テクニカルリードとして、これら数理的知見を実務の動態解析へと昇華させることは、経験則を超えた「予測可能なシステム設計」を実現するための確固たる基盤となるのである。
\end{enumerate}

\end{document}
